%% toc: P2 15 15.1,15.2,15.3,15.8
\section{Decision bounds, reductions and checked certificates for propositional satisfiability}
\label{sec:p02-w02-p02-018-decision-certificates}

\emph{Purpose.}  The propositional resolution rule, its soundness and its
refutation completeness are the exact inherited results of
\cref{sec:prop-metatheory-04}, with the refutation displayed as a graph in
Figure~\ref{fig:p02-resolution-dag}.  This section supplies what those results
do not: an explicit input representation for the four finite decision problems
of classical propositional logic, a truth-table procedure for them with a
written termination, correctness and running-time proof, the exact limits of
that decidability claim, the proofs that \(\mathsf{SAT}\) is in
\(\mathsf{NP}\), that \(\mathsf{TAUT}\) is in \(\mathsf{coNP}\), that negation
interreduces \(\mathsf{SAT}\) with the complement of \(\mathsf{TAUT}\), and
that every \(\mathsf{NP}\) problem reduces to \(\mathsf{SAT}\), and finally
three checked certificate formats with a worked instance of each, together
with a fully annotated first-order resolution refutation.

\emph{Inputs.}  Finite strings over the fixed token alphabet of
\cref{def:p02-w02-p02-018-input}, presenting a formula, a finite list of
premises with a conclusion, or a pair of formulas; and, for the certificate
checks, finite clause sets with a candidate valuation or a candidate
refutation log.  \emph{Conclusions.}  A yes or no answer with a witness
valuation where one exists, a bound on the number of machine steps, and
checked evidence for satisfiability, unsatisfiability and failure of
consequence.

\emph{Standing notation.}  Formulas are the members of
\(\mathsf{Fm}_{\mathsf{CPL}}\) of \cref{sec:prop-syntax-semantics-01}, with
primitive \(\bottom,\neg,\land,\lor,\to\) and variables \(p_n\)
(\(n\in\N\)).  A valuation \(v\) assigns \(\mathbf T\) or \(\mathbf F\) to
every variable, \(\widehat v\) is its extension to formulas, and
satisfaction, satisfiability, validity, semantic consequence
\(\Gamma\entails\varphi\) and logical equivalence are those of
\cref{sec:prop-syntax-semantics-02}.  \(\operatorname{Var}(e)\) is the finite
set of variables occurring in \(e\).  Clauses, the clause interpretation
\(\operatorname{form}_{\mathsf{CPL}}\) and literal complements
\(\overline L\) are those of \cref{sec:prop-syntax-semantics-04,%
sec:prop-metatheory-04}.  In examples \(p,q,r\) are the distinct variables
\(p_0,p_1,p_2\).  The complexity classes, reductions and the languages
\(\mathsf{SAT}\) and \(\mathsf{TAUT}\) are those of
\cref{def:p02-complexity-standing,def:p02-sat-taut}.

\subsection{The declared input representation}
\label{sec:p02-w02-p02-018-input}

\begin{hbconvention}{def:p02-w02-p02-018-input}{Effective finite inputs for classical propositional logic}
\emph{Tokens.}  The token alphabet consists of the thirteen distinct symbols
\[
 \texttt{p},\ \texttt{0},\ \texttt{1},\ \texttt{;},\ \bottom,\ \neg,\
 \land,\ \lor,\ \to,\ (,\ ),\ \texttt{,}\ ,\ \texttt{|}\ ,
\]
listed here in a fixed order; the token in position \(i\) of this list
(\(0\leq i\leq 12\)) is encoded by the four-bit binary numeral of \(i\).
Every input is a bit string; it is \emph{decodable} when its length is a
multiple of \(4\) and each consecutive group of four bits is the code of a token.
Write \(m\) for the bit length of an input, so a decodable input has \(m/4\)
tokens.

\emph{Formulas.}  The variable \(p_n\) is written as the token \texttt{p},
followed by the binary numeral of \(n\) without leading zeros (the single
digit \texttt{0} when \(n=0\)), followed by the terminator \texttt{;}.  A
formula is written as its fully parenthesized string of
\cref{sec:prop-syntax-semantics-01} with each variable so written; the
separators \texttt{,} and \texttt{|} do not occur in formula strings.  This
is the formula encoding of \cref{def:p02-sat-taut} with its code table
fixed.

\emph{Instances.}  The four decision problems are the following sets of bit
strings.
\begin{enumerate}
\item \(\mathsf{SAT}\): encodings of satisfiable formulas.
\item \(\mathsf{TAUT}\): encodings of valid formulas.
\item \(\mathsf{CONS}\): encodings of strings
      \(\gamma_1\,\texttt{,}\,\gamma_2\,\texttt{,}\cdots\texttt{,}\,\gamma_r\,
      \texttt{|}\,\varphi\) with \(r\geq0\) formula strings before the
      separator \texttt{|} (for \(r=0\) the string begins with \texttt{|})
      such that \(\{\gamma_1,\ldots,\gamma_r\}\entails\varphi\).  The list
      presents the finite premise set \(\Gamma=\{\gamma_1,\ldots,\gamma_r\}\);
      order and repetition in the list do not affect membership.
\item \(\mathsf{EQV}\): encodings of strings
      \(\varphi\,\texttt{|}\,\psi\) of two formula strings such that
      \(\varphi\) and \(\psi\) are logically equivalent.
\end{enumerate}
A bit string that is not decodable, or whose token string does not have the
displayed shape with formula strings in the displayed positions, is
\emph{malformed} and belongs to none of the four sets.  For an instance
\(x\), \(k\) denotes the number of distinct variables occurring in it.
Because every variable occurrence uses at least three tokens,
\(k\leq m/12\).

The names \(\mathsf{CONS}\) and \(\mathsf{EQV}\) are needed to state the
consequence and equivalence problems as sets of strings beside the two
languages of \cref{def:p02-sat-taut}; they are metalinguistic names of
languages, not connectives, and are used in
\cref{thm:p02-w02-p02-018-decision}.  Any other representation reaches these
results only through a computable translation into this one.
\index{CONS@\(\mathsf{CONS}\)}\index{EQV@\(\mathsf{EQV}\)}
\index{input representation!propositional decision problems}
\end{hbconvention}

\subsection{The truth-table decision procedure and its bound}
\label{sec:p02-w02-p02-018-truth-table}

\begin{hbmethod}{def:p02-w02-p02-018-evaluator}{Stack evaluation of formula strings}
\emph{Input.}  A token string \(w\) without separators, and a
\emph{valuation table} \(\tau\): a finite list of pairs \((n,b)\), each a
binary numeral \(n\) and a truth value \(b\), with no numeral listed twice.
The table determines the valuation \(v_\tau\) that gives \(p_n\) the value
\(b\) when \((n,b)\) is listed and \(\mathbf F\) otherwise.

\emph{Procedure.}  Start with an empty stack whose entries are truth values
or the markers \(\neg\), \((\), \(\land\), \(\lor\), \(\to\).  To
\emph{reduce} the stack, repeat: while the top entry is a truth value \(b\)
and the entry below it is \(\neg\), replace these two entries by the
opposite value of \(b\).  Read \(w\) from left to right.
\begin{enumerate}
\item On \texttt{p}: read the following digit tokens up to the next
      \texttt{;}.  If no \texttt{;} follows, if the digit string is empty, or
      if it has a leading \texttt{0} and more than one digit, stop with the
      answer \emph{malformed}.  Otherwise push the value that \(v_\tau\)
      gives to the variable with this numeral, then reduce.
\item On \(\bottom\): push \(\mathbf F\), then reduce.
\item On \(\neg\), \((\), \(\land\), \(\lor\) or \(\to\): push it.
\item On \()\): if the four top entries are, from the top, a truth value
      \(b\), a binary connective \(\circ\), a truth value \(a\) and \((\),
      replace them by the value of \(a\circ b\) under the truth clause of
      \(\circ\), then reduce; otherwise stop with \emph{malformed}.
\item On any other token stop with \emph{malformed}.
\end{enumerate}
At the end of \(w\), answer \emph{well formed with value \(b\)} if the stack
consists of exactly one truth value \(b\), and \emph{malformed} otherwise.
\end{hbmethod}

\begin{hblemma}{thm:p02-w02-p02-018-evaluator}{proved=annex:p02-w02-p02-018-evaluator}{Correctness and cost of stack evaluation}
Let \(w\) be a token string without separators and \(\tau\) a valuation
table.
\begin{enumerate}
\item Method~\ref{def:p02-w02-p02-018-evaluator} answers \emph{well formed}
      on \(w\) if and only if \(w\) is the string of a formula \(\chi\), and
      then its value is \(\widehat{v_\tau}(\chi)\).
\item On a deterministic machine with a read-only input tape holding \(w\)
      and work tapes holding \(\tau\) and the stack, the method halts within
      \(c\,t\,(t+s)\) steps, where \(t\) is the number of tokens of \(w\),
      \(s\) the number of symbols of the written table and \(c\) a constant
      independent of \(w\) and \(\tau\).
\end{enumerate}
\end{hblemma}

\begin{hbmethod}{def:p02-w02-p02-018-truth-table}{Truth-table decision procedure}
\emph{Input.}  A bit string \(x\) and one of the four problems of
\cref{def:p02-w02-p02-018-input}.

\emph{Step 1 (recognition).}  Decode \(x\) into tokens; split the token string
at its separators; check that the separators have the shape required by the
problem; and run Method~\ref{def:p02-w02-p02-018-evaluator} with the empty
table on each segment.  If any check fails, \emph{reject}.

\emph{Step 2 (variables).}  List the distinct variables of the instance as
\(x_1,\ldots,x_k\) in order of first occurrence, and write the table
\(\tau=((x_1,\mathbf F),\ldots,(x_k,\mathbf F))\).

\emph{Step 3 (rows).}  Visit the \(2^k\) assignments
\(u:\{x_1,\ldots,x_k\}\to\{\mathbf T,\mathbf F\}\) in lexicographic order of
\((u(x_1),\ldots,u(x_k))\) with \(\mathbf F\) before \(\mathbf T\); the table
holds the current row, and the next row is obtained by changing the final
run of \(\mathbf T\) entries to \(\mathbf F\) and the entry before that
run from \(\mathbf F\) to \(\mathbf T\).  For each row evaluate every segment
under \(v_\tau\) by Method~\ref{def:p02-w02-p02-018-evaluator} and stop at the
first row on which
\begin{itemize}
\item for \(\mathsf{SAT}\), the formula is true: \emph{accept} and output the
      row;
\item for \(\mathsf{TAUT}\), the formula is false: \emph{reject} and output
      the row;
\item for \(\mathsf{CONS}\), every premise is true and the conclusion is
      false: \emph{reject} and output the row;
\item for \(\mathsf{EQV}\), the two formulas receive different values:
      \emph{reject} and output the row.
\end{itemize}
\emph{Step 4 (exhaustion).}  If no row stops the search, \emph{reject} for
\(\mathsf{SAT}\) and \emph{accept} for \(\mathsf{TAUT}\), \(\mathsf{CONS}\)
and \(\mathsf{EQV}\).
\index{truth-table decision procedure}
\end{hbmethod}

\begin{hbtheorem}{thm:p02-w02-p02-018-decision}{proved=annex:p02-w02-p02-018-decision}{Decidability with a running-time bound}
For each of \(\mathsf{SAT}\), \(\mathsf{TAUT}\), \(\mathsf{CONS}\) and
\(\mathsf{EQV}\), Method~\ref{def:p02-w02-p02-018-truth-table} halts on every
bit string \(x\) and accepts \(x\) exactly when \(x\) belongs to the problem.
A row output on acceptance for \(\mathsf{SAT}\), extended by \(\mathbf F\) on
all other variables, satisfies the formula; a row output on rejection for
\(\mathsf{TAUT}\), \(\mathsf{CONS}\) or \(\mathsf{EQV}\), so extended,
falsifies the formula, is a countervaluation to
\(\Gamma\entails\varphi\), or separates the two formulas.  On a deterministic
multitape machine the procedure halts within
\[
 c\,(2^{k}+1)\,m^{2}
\]
steps, where \(m\) is the bit length of \(x\), \(k\) the number of distinct
variables of the instance (\(k=0\) for a malformed input) and \(c\) a
constant independent of \(x\).  Every valid formula with \(k\) variables
makes the procedure visit all \(2^{k}\) rows.
\end{hbtheorem}

The bound has two factors: the number of rows, \(2^{k}\), and the cost of one
row, which is at most quadratic in the input length.  The last sentence of
the theorem shows that this procedure attains the exponential factor.  That
is not a lower bound for the problem: a proof that no polynomial-time
procedure decides \(\mathsf{SAT}\) would settle \(\mathsf P\ne\mathsf{NP}\),
since \(\mathsf{SAT}\) is \(\mathsf{NP}\)-complete
(\cref{thm:p02-w02-p02-018-cook-levin}).  The theorem also gives a written
proof of the content of the cited \cref{thm:p02-decidability} for the
representation of \cref{def:p02-w02-p02-018-input}.

\begin{hbwarning}{def:p02-w02-p02-018-scope}{Limits of the decidability theorem}
\Cref{thm:p02-w02-p02-018-decision} decides consequence from a
\emph{finite} premise list presented by the declared encoding, over the
countable variable set \(\{p_n:n\in\N\}\).  It gives no procedure for a fixed
infinite theory \(\Gamma\), and none exists in general: by
\cref{thm:p02-w02-p02-018-limits}(1) there are sets
\(\Gamma\subseteq\{p_n:n\in\N\}\) whose consequence sets are not decided by
any machine.  Compactness (\cref{thm:prop-compactness}) implies that a
consequence of an infinite \(\Gamma\) is already a consequence of a finite
subset, but it does not say how to find that subset from \(\Gamma\).  Nor does the theorem apply to the
languages \(\operatorname{Form}(P)\) of
\cref{def:p02-arbitrary-prop-language} over an arbitrary set \(P\): a
decision procedure acts on finite strings, and by
\cref{thm:p02-w02-p02-018-limits}(2) an uncountable \(P\) has no injective
representation by finite strings at all.  A countable \(P\) is covered only
when it is presented by a computable translation into the declared
encoding.
\end{hbwarning}

\begin{hbtheorem}{thm:p02-w02-p02-018-limits}{proved=annex:p02-w02-p02-018-limits}{Infinite theories and uncountable alphabets}
\begin{enumerate}
\item For \(A\subseteq\N\) let \(\Gamma_A=\{p_n:n\in A\}\).  Then
      \(\Gamma_A\entails p_n\) if and only if \(n\in A\).  Consequently
      the map \(A\mapsto\{\varphi:\Gamma_A\entails\varphi\}\) is injective,
      and for all but countably many \(A\subseteq\N\) the set of encodings
      of formulas \(\varphi\) with \(\Gamma_A\entails\varphi\) is decided by no
      deterministic multitape machine.
\item If \(P\) is uncountable, there is no injective map from \(P\) into the
      set \(\{0,1\}^*\) of finite bit strings.
\end{enumerate}
\end{hbtheorem}

\subsection{Membership, the negation reduction and \texorpdfstring{\(\mathsf{NP}\)}{NP}-hardness}
\label{sec:p02-w02-p02-018-complexity}

\begin{hbconvention}{def:p02-w02-p02-018-verifier}{Verifiers and pairing}
For bit strings \(x=x_1\cdots x_n\) and \(y\), put
\(\langle x,y\rangle\defeq x_1x_1x_2x_2\cdots x_nx_n\,01\,y\).  Reading
the string in consecutive pairs of bits from the left, the first pair
\(01\) separates the two parts, because every earlier pair is \(00\) or
\(11\); so the pairing is injective and both parts are recovered in one
scan.  A predicate
\(R\) of two strings is \emph{polynomial time} when a deterministic multitape
machine decides \(\{\langle x,y\rangle:R(x,y)\}\) in a number of steps bounded
by a polynomial in the length of \(\langle x,y\rangle\).  This fixes the
predicate clause of \cref{def:p02-complexity-standing}.  For
\(L\subseteq\{0,1\}^*\), \(L^{\mathsf c}\) is its complement
\(\{0,1\}^*\setminus L\); thus \(\mathsf{TAUT}^{\mathsf c}\) consists of the
malformed strings and the encodings of formulas that are not valid.
\end{hbconvention}

\begin{hbtheorem}{thm:p02-w02-p02-018-membership}{proved=annex:p02-w02-p02-018-membership}{\(\mathsf{SAT}\) is in \(\mathsf{NP}\) and \(\mathsf{TAUT}\) is in \(\mathsf{coNP}\)}
\(\mathsf{SAT}\in\mathsf{NP}\), witnessed by the certificate that lists the
values of a satisfying valuation on the variables of the formula in order of
first occurrence, checked by one evaluation.  \(\mathsf{TAUT}^{\mathsf c}\in
\mathsf{NP}\), witnessed by a falsifying valuation listed in the same way (or
the empty string for a malformed input); hence \(\mathsf{TAUT}\in
\mathsf{coNP}\).  In both cases the certificate has length at most \(m\) and
the check takes at most \(c\,m^{2}\) steps.
\end{hbtheorem}

\begin{hbtheorem}{thm:p02-w02-p02-018-negation}{proved=annex:p02-w02-p02-018-negation}{The negation reduction}
Let \(h(x)\) be the encoding of \(\neg\varphi\) when \(x\) encodes the formula
\(\varphi\), and the encoding of \(\neg\bottom\) when \(x\) is malformed.
Then \(h\) is computable within \(c\,(m+1)^{2}\) steps on inputs of bit
length \(m\), for a constant \(c\), and, for every bit string \(x\),
\[
 x\in\mathsf{SAT}\iff h(x)\in\mathsf{TAUT}^{\mathsf c},
 \qquad
 x\in\mathsf{TAUT}^{\mathsf c}\iff h(x)\in\mathsf{SAT}.
\]
Consequently, for every \(L\subseteq\{0,1\}^*\), \(L\) has a polynomial-time
many-one reduction to \(\mathsf{SAT}\) if and only if \(L^{\mathsf c}\) has one
to \(\mathsf{TAUT}\).
\end{hbtheorem}

\begin{hbtheorem}{thm:p02-w02-p02-018-cook-levin}{proved=annex:p02-w02-p02-018-cook-levin}{Every \(\mathsf{NP}\) problem reduces to \(\mathsf{SAT}\)}
For every \(L\in\mathsf{NP}\) there is a polynomial-time computable map
\(x\mapsto\varphi_x\), with \(\varphi_x\) an encoded formula, such that
\(x\in L\) iff \(\varphi_x\) is satisfiable.  With
\cref{thm:p02-w02-p02-018-membership}, \(\mathsf{SAT}\) is
\(\mathsf{NP}\)-complete under polynomial-time many-one reductions.
\end{hbtheorem}

\begin{hbcorollary}{thm:p02-w02-p02-018-taut-conp}{proved=annex:p02-w02-p02-018-taut-conp}{\(\mathsf{TAUT}\) is \(\mathsf{coNP}\)-complete}
\(\mathsf{TAUT}\) is \(\mathsf{coNP}\)-complete under polynomial-time
many-one reductions.
\end{hbcorollary}

The formula \(\varphi_x\) describes a run of a polynomial-time verifier for
\(L\) on the input \(x\) and an unknown certificate, step by step, in
variables that record the machine state, the head positions and the cell
contents at each time.  This tableau-of-computation method is the
construction of \textcite[p.~152, Theorem~1]{Cook1971}, which encodes a
one-tape nondeterministic machine and states the result for tautologies in
disjunctive normal form under polynomial-time oracle reducibility.  The
annex encodes a deterministic multitape verifier instead, so that the
reduction is many-one and uses the machine model of
\cref{def:p02-complexity-standing} without a tape-reduction step.  The
reconstruction in \cref{sec:p09-presentation-n-c-026} encodes a one-tape
verifier by windows of three adjacent cells and so presupposes a one-tape
simulation of the verifier; the encoding here follows every head of the
multitape verifier with its own position variables instead.
\Cref{thm:p02-w02-p02-018-cook-levin,thm:p02-w02-p02-018-taut-conp} prove the
statements cited as \cref{thm:p02-sat-np-complete,thm:p02-taut-conp-complete}
and \cref{thm:p04-cook-levin}.  They classify the problems; they do not decide
whether \(\mathsf P=\mathsf{NP}\) or \(\mathsf{NP}=\mathsf{coNP}\).

\subsection{Checked certificates}
\label{sec:p02-w02-p02-018-certificates}

\begin{hbdefinition}{def:p02-w02-p02-018-formats}{Three certificate formats and their checks}
Let \(F\) and \(S\) be finite clause sets, \(\Gamma\) a finite formula set and
\(\varphi\) a formula.
\begin{enumerate}
\item A \emph{satisfying-valuation certificate} for \(F\) is a total
      assignment of truth values to \(\operatorname{Var}(F)\).  Its
      \emph{clause check} visits each clause \(C\in F\) and looks for a literal
      of \(C\) made true by the assignment; it accepts when every clause has
      one.  This is the SAT certificate of
      \cref{def:p12-sat-certificate-boundary}.
\item A \emph{resolution refutation log} for \(S\) is a finite list of
      numbered lines \(1,\ldots,N\).  Line \(i\) carries a clause \(C_i\)
      and a justification, which is either \emph{input} or a triple
      \((j,k,L)\) with \(j,k<i\): the parents \(C_j\), \(C_k\) and the
      pivot literal \(L\).  The \emph{line check} accepts an input line when
      \(C_i\in S\), and a triple line when \(L\in C_j\),
      \(\overline L\in C_k\) and
      \(C_i=(C_j\setminus\{L\})\cup(C_k\setminus\{\overline L\})\).  The
      \emph{log check} accepts when every line passes and \(C_N=\varnothing\).
\item A \emph{countervaluation certificate} for \(\Gamma\entails\varphi\) is
      a total assignment of truth values to
      \(\operatorname{Var}(\Gamma\cup\{\varphi\})\).  Its check evaluates every
      member of \(\Gamma\) and \(\varphi\) under the assignment, extended by
      \(\mathbf F\) elsewhere, and accepts when every member of \(\Gamma\) is
      true and \(\varphi\) is false.
\end{enumerate}
These are finite evidence objects with mechanical checks, distinct from the
semantic facts they establish; the distinction is used in
\cref{thm:p02-w02-p02-018-checks}.
\index{certificate!satisfying valuation}\index{resolution refutation log}
\index{certificate!countervaluation}
\end{hbdefinition}

\begin{hbtheorem}{thm:p02-w02-p02-018-checks}{proved=annex:p02-w02-p02-018-checks}{Soundness and completeness of the three checks}
\begin{enumerate}
\item The clause check accepts an assignment for \(F\) if and only if every
      valuation extending it satisfies \(F\).  So \(F\) is satisfiable iff
      some certificate is accepted.
\item If the log check accepts a log for \(S\), then \(S\) is not
      satisfiable.  If \(S\) is not satisfiable, some log for \(S\) is
      accepted.
\item The countervaluation check accepts an assignment if and only if every
      valuation extending it satisfies \(\Gamma\) and falsifies \(\varphi\).
      So \(\Gamma\not\entails\varphi\) iff some certificate is accepted.
\end{enumerate}
\end{hbtheorem}

\begin{hblemma}{thm:p02-w02-p02-018-extraction}{proved=annex:p02-w02-p02-018-extraction}{Countervaluations from failed searches}
Let \(\Gamma\) be finite.
\begin{enumerate}
\item If Method~\ref{def:p02-w02-p02-018-truth-table} rejects an instance of
      \(\mathsf{CONS}\) presenting \(\Gamma\) and \(\varphi\), its output row,
      restricted to \(\operatorname{Var}(\Gamma\cup\{\varphi\})\), is accepted
      by the countervaluation check.
\item Let \(F_\Gamma\) and \(F_\neg\) be finite clause sets such that every
      valuation satisfying \(F_\Gamma\) satisfies \(\Gamma\), every valuation
      satisfying \(F_\neg\) falsifies \(\varphi\), and
      \(\operatorname{Var}(\Gamma\cup\{\varphi\})\subseteq
      \operatorname{Var}(F_\Gamma\cup F_\neg)\).  If
      Method~\ref{def:p12-dpll} returns a valuation on
      \(F_\Gamma\cup F_\neg\), its restriction to
      \(\operatorname{Var}(\Gamma\cup\{\varphi\})\) is accepted by the
      countervaluation check.
\end{enumerate}
\end{hblemma}

The hypotheses of the second clause hold when each member of \(\Gamma\) and
the formula \(\neg\varphi\) is the conjunctive normal form, in the same
variables, of the corresponding clause set
(\cref{lem:prop-clause-satisfaction}).  They fail for a Tseitin clause set,
which is only equisatisfiable: there the returned valuation must first be
restricted, and its correctness rests on the definitional clauses, not on this
lemma.  A finished tableau with an open branch gives the third route, by
\cref{thm:p02-w02-p02-014-consequence}.

\begin{hbexample}{def:p02-w02-p02-018-sat-example}{A checked satisfying valuation}
Let
\[
 F=\bigl\{\{p,q\},\{\neg p,r\},\{\neg q,r\}\bigr\},
\]
so \(\operatorname{Var}(F)=\{p,q,r\}\).  The certificate is the total assignment
\(v(p)=\mathbf F\), \(v(q)=\mathbf T\), \(v(r)=\mathbf T\).  The clause check
runs as follows.
\[
\begin{array}{lll}
\text{clause} & \text{literal found} & \text{value}\\\hline
\{p,q\} & q & \mathbf T\\
\{\neg p,r\} & \neg p & \mathbf T\\
\{\neg q,r\} & r & \mathbf T
\end{array}
\]
Every clause has a true literal, so the check accepts and \(F\) is satisfiable
by \cref{thm:p02-w02-p02-018-checks}(1).  The candidate
\(p,q,r\mapsto\mathbf F\) is rejected at the first clause, which has no true
literal; a rejected candidate shows nothing about \(F\).
\end{hbexample}

\begin{hbexample}{def:p02-w02-p02-018-log-example}{A checked resolution refutation log}
Let \(S=F\cup\bigl\{\{\neg r\}\bigr\}\) with \(F\) as in
\cref{def:p02-w02-p02-018-sat-example}.  This is the clause set of
Figure~\ref{fig:p02-resolution-dag} and of the trace
\cref{def:p12-dpll-trace}.  The log and its line checks are:
\[
\begin{array}{rlll}
i & C_i & \text{justification} & \text{line check}\\\hline
1 & \{p,q\} & \text{input} & \in S\\
2 & \{\neg p,r\} & \text{input} & \in S\\
3 & \{\neg q,r\} & \text{input} & \in S\\
4 & \{\neg r\} & \text{input} & \in S\\
5 & \{\neg p\} & (2,4,r) & r\in C_2,\ \neg r\in C_4,\
   \{\neg p\}\cup\varnothing=\{\neg p\}\\
6 & \{\neg q\} & (3,4,r) & r\in C_3,\ \neg r\in C_4,\
   \{\neg q\}\cup\varnothing=\{\neg q\}\\
7 & \{q\} & (1,5,p) & p\in C_1,\ \neg p\in C_5,\
   \{q\}\cup\varnothing=\{q\}\\
8 & \varnothing & (7,6,q) & q\in C_7,\ \neg q\in C_6,\
   \varnothing\cup\varnothing=\varnothing
\end{array}
\]
Every line passes and the last clause is empty, so the log check accepts and
\(S\) is not satisfiable by \cref{thm:p02-w02-p02-018-checks}(2).  Since
\(\neg r=\operatorname{form}_{\mathsf{CPL}}(\{\neg r\})\), the same log
certifies the consequence
\(\{p\lor q,\ \neg p\lor r,\ \neg q\lor r\}\entails r\): a valuation
satisfying the three premises and falsifying \(r\) would satisfy \(S\).
\end{hbexample}

\begin{hbexample}{def:p02-w02-p02-018-counter-example}{A countervaluation extracted from a failed search and checked}
Let \(\Gamma=\{p\lor q,\ \neg p\lor r\}\) and \(\varphi=r\).  With
\(p,q,r=p_0,p_1,p_2\), the members of \(\Gamma\) are
\(\operatorname{form}_{\mathsf{CPL}}(\{p,q\})\) and
\(\operatorname{form}_{\mathsf{CPL}}(\{\neg p,r\})\), and \(\neg\varphi\)
is satisfied exactly by the valuations satisfying \(\{\neg r\}\).  So the
hypotheses of \cref{thm:p02-w02-p02-018-extraction}(2) hold for
\(F_\Gamma=\bigl\{\{p,q\},\{\neg p,r\}\bigr\}\) and
\(F_\neg=\bigl\{\{\neg r\}\bigr\}\).

\emph{Search.}  Method~\ref{def:p12-dpll} starts from
\(\bigl\{\{p,q\},\{\neg p,r\},\{\neg r\}\bigr\}\) with the empty assignment.
The unit clause \(\{\neg r\}\) sets \(r=\mathbf F\), leaving
\(\bigl\{\{p,q\},\{\neg p\}\bigr\}\).  The unit clause \(\{\neg p\}\) sets
\(p=\mathbf F\), leaving \(\bigl\{\{q\}\bigr\}\).  The unit clause \(\{q\}\)
sets \(q=\mathbf T\), leaving no clause.  The search returns
\(v(p)=\mathbf F\), \(v(q)=\mathbf T\), \(v(r)=\mathbf F\), with no
unassigned variable left to extend.

\emph{Check.}  \(\widehat v(p\lor q)=\mathbf T\) because
\(v(q)=\mathbf T\); \(\widehat v(\neg p\lor r)=\mathbf T\) because
\(v(p)=\mathbf F\); and \(\widehat v(r)=\mathbf F\).  The countervaluation
check accepts, so \(\Gamma\not\entails r\) by
\cref{thm:p02-w02-p02-018-checks}(3).

\emph{The same answer by truth table.}  Method~\ref{def:p02-w02-p02-018-truth-table}
on the instance \(p\lor q\,\texttt{,}\,\neg p\lor r\,\texttt{|}\,r\) lists
\(x_1,x_2,x_3=p,q,r\) and visits \(\mathbf{FFF}\), \(\mathbf{FFT}\),
\(\mathbf{FTF}\).  The first row falsifies \(p\lor q\), the second makes
\(r\) true, and the third is the row found above, so the procedure rejects at
the third of the \(2^3\) rows and outputs it.  Adding the premise
\(\neg q\lor r\) removes this countervaluation; the log of
\cref{def:p02-w02-p02-018-log-example} then certifies the consequence.
\end{hbexample}

\begin{hbwarning}{def:p02-w02-p02-018-bad-log}{A line that resolves on two pivots is rejected}
From \(\{p,q\}\) and \(\{\neg p,\neg q\}\) a line claiming
\(\varnothing\) fails the line check: on the pivot \(p\) the resolvent is
\(\{q,\neg q\}\), on the pivot \(q\) it is \(\{p,\neg p\}\), and neither is
empty.  The rule removes exactly one complementary pair.  Accepting such a line
would be unsound, since \(p\mapsto\mathbf T\), \(q\mapsto\mathbf F\) satisfies
both clauses.  Likewise an assignment that leaves a variable of \(F\)
unassigned is not a certificate in the sense of
\cref{def:p02-w02-p02-018-formats}(1), and an accepted countervaluation for
\(\Gamma\entails\varphi\) shows only that this consequence fails, not that
\(\Gamma\entails\neg\varphi\).
\end{hbwarning}

\subsection{A fully annotated first-order resolution refutation}
\label{sec:p02-w02-p02-018-fol-dag}

The first-order clauses, substitutions, renaming apart, most general unifiers
and the resolution and factoring rules are those of
\cref{sec:fol-cardinality-normalforms-04}.  In
Figure~\ref{fig:p02-w02-p02-018-fol-dag} every inference is an instance of
the first-order resolution rule with the parents
\(E\ni A\) and \(F\ni\neg B\); each edge pair is annotated with the rule, the
resolved literal pair \(A\) and \(\neg B\), and the most general unifier
\(\theta\) of \(A\) and \(B\).  For a ground step the unifier is the identity
\(\mathrm{id}\), the substitution with empty support.  No factoring step is
used, and no renaming is needed: the parents of the first step have the
disjoint variable sets \(\{x\}\) and \(\{y\}\), and every later step has a
ground parent.

\begin{figure}[htbp]
\centering
\begin{tikzpicture}[
  font=\small,
  clause/.style={draw,rounded corners,align=center,inner sep=3pt},
  note/.style={align=left,font=\footnotesize,inner sep=2pt},
  >=Stealth]
  \node[clause] (c1) at (0,0) {\(C_1=\{P(x),Q(x)\}\)};
  \node[clause] (c2) at (4.6,0) {\(C_2=\{\neg Q(y),R(y)\}\)};
  \node[clause] (c3) at (9.0,-1.9) {\(C_3=\{\neg P(a)\}\)};
  \node[clause] (c4) at (9.0,-3.8) {\(C_4=\{\neg R(a)\}\)};
  \node[clause] (c5) at (2.3,-1.9) {\(C_5=\{P(y),R(y)\}\)};
  \node[clause] (c6) at (4.6,-3.8) {\(C_6=\{R(a)\}\)};
  \node[clause] (c7) at (6.8,-5.7) {\(C_7=\varnothing\)};
  \draw[->] (c1) -- (c5);
  \draw[->] (c2) -- (c5);
  \draw[->] (c5) -- (c6);
  \draw[->] (c3) -- (c6);
  \draw[->] (c6) -- (c7);
  \draw[->] (c4) -- (c7);
  \node[note,anchor=east] at (1.6,-1.9)
    {Resolution\\\(A=Q(x)\in C_1\)\\\(\neg B=\neg Q(y)\in C_2\)\\\(\theta_1=\{x\mapsto y\}\)};
  \node[note,anchor=east] at (3.4,-3.8)
    {Resolution\\\(A=P(y)\in C_5\)\\\(\neg B=\neg P(a)\in C_3\)\\\(\theta_2=\{y\mapsto a\}\)};
  \node[note,anchor=east] at (5.6,-5.7)
    {Resolution\\\(A=R(a)\in C_6\)\\\(\neg B=\neg R(a)\in C_4\)\\\(\theta_3=\mathrm{id}\) (ground)};
\end{tikzpicture}
\caption{A first-order resolution refutation from the input clauses
\(C_1,\ldots,C_4\) to the empty clause \(C_7\).  Each derived clause is the
resolvent of its two incoming parents: \(C_5\) of \(C_1,C_2\); \(C_6\) of
\(C_5,C_3\); \(C_7\) of \(C_6,C_4\).  Each annotation names the rule, the
resolved literal pair and the most general unifier used; the last step is
ground and uses the identity substitution.}
\label{fig:p02-w02-p02-018-fol-dag}
\end{figure}
\index{resolution!annotated refutation DAG}

\begin{hblemma}{thm:p02-w02-p02-018-fol-dag}{proved=annex:p02-w02-p02-018-fol-dag}{The annotated refutation is correct}
In Figure~\ref{fig:p02-w02-p02-018-fol-dag}, each \(\theta_i\) is a most
general unifier of the displayed atoms \(A\) and \(B\), and is the output of
the unification procedure of \cref{sec:fol-cardinality-normalforms-04} on
them; each derived clause is the displayed resolvent
\(\bigl((E\setminus\{A\})\cup(F\setminus\{\neg B\})\bigr)\theta_i\).
Consequently \(\{C_1,C_2,C_3,C_4\}\) is not satisfiable.
\end{hblemma}

Figure~\ref{fig:p02-fol-resolution-dag} shows a two-step refutation whose
second step, a ground resolution of \(Q(a)\) against \(\neg Q(a)\), uses the
identity unifier; Figure~\ref{fig:p02-w02-p02-018-fol-dag} states every such
annotation explicitly.

\subsection{Relationships and sources}
\label{sec:p02-w02-p02-018-relationships}

The truth-table procedure is one of four decision routes for the same finite
problems.  Resolution decides unsatisfiability of a finite clause set by
saturation, by \cref{thm:prop-resolution-complete}; the systematic tableau
procedure decides satisfiability, validity and finite-premise consequence by
\cref{thm:p02-w02-p02-014-decision}; and the search of
Method~\ref{def:p12-dpll} decides satisfiability of a finite clause set by
\cref{thm:p12-dpll-correctness}.  Only the truth-table bound is stated in
terms of the variable count alone; the polynomial special cases of
Table~\ref{tab:p02-tractable-sat} restrict the input syntax instead.  The
class and reduction conventions are those of
\cref{def:p02-complexity-standing}, and Part~4 places the classifications of
\cref{sec:p02-w02-p02-018-complexity} in its class ladder
(\cref{def:p04-complexity-convention}).  The resolution method is due to
\textcite{Robinson1965}; the reduction of
\cref{thm:p02-w02-p02-018-cook-levin} follows the method of
\textcite[p.~152, Theorem~1]{Cook1971}.
